% F: registro real (trace do Langfuse) de uma execução nas duas abordagens com agente, pergunta a01, Claude Haiku 5.5.
% Consulta, argumentos e saídas copiados dos traces da rodada final (final_results.jsonl, campo trace_url);
% trechos longos abreviados com [...]. As duas abordagens acertaram a a01 nos quatro modelos.
% Sem tabular: só parágrafos com recuo e filetes, para sair igual em qualquer versão do TeX Live (com o pacote
% array de 2026, células p{} com caixas internas sobrepunham as linhas no Overleaf).
% Usa \ttfamily e \char34 (e não verbatim) porque o babel torna a aspa dupla ativa e o verbatim quebra acentos.
\providecommand{\jq}{\char34{}}
\newlength{\traceind}\setlength{\traceind}{0.21\columnwidth}
\newcommand{\tracerule}[1]{\par\vspace{2pt}\noindent\rule{\columnwidth}{#1}\par\vspace{3pt}}
\newcommand{\tracehead}[1]{\par\noindent\textbf{#1}\par\vspace{2pt}}
\newcommand{\tracestep}[2]{\par\noindent\hangindent=\traceind\hangafter=1\makebox[\traceind][l]{\itshape #1}#2\par}
\newcommand{\tracecode}[1]{{\raggedright\leftskip=\traceind\ttfamily\scriptsize\frenchspacing\vspace{2pt}#1\par}\vspace{3pt}}
\newcommand{\traceloop}{{\leftskip=\traceind\scriptsize\itshape $\uparrow$ volta ao modelo; pode chamar de novo
  (até \nMaxCalls{})\par}\vspace{2pt}}
\vspace{2mm}
\noindent\begin{minipage}{\columnwidth}
\footnotesize\raggedright\setlength{\parindent}{0pt}
\tracerule{0.8pt}
Pergunta (a01): Quem deu mais passes certos no terço final do campo pela Argentina?
\tracerule{0.5pt}
\tracehead{\armTextToCypher{}}
\tracestep{Modelo}{chama \texttt{run\_cypher} com a consulta que escreveu:}
\tracecode{MATCH (j:Jogador)-[r:REALIZOU]->\\
\ \ (p:Partida [...])\\
WHERE j.time = 'Argentina'\\
\ \ AND r.acao = 'passe' AND r.sucesso = true\\
\ \ AND r.terco = 'ataque'\\
RETURN j.nome AS jogador,\\
\ \ count(r) AS passes\_certos\_terco\_final\\
ORDER BY passes\_certos\_terco\_final DESC\\
LIMIT 10}
\tracestep{Ferramenta}{executa essa consulta no grafo e devolve:}
\tracecode{Lionel Andrés Messi Cuccittini: 24\\ Alexis Mac Allister: 18 [...]}
\traceloop
\tracestep{Modelo}{responde: Lionel Andrés Messi Cuccittini}
\tracerule{0.5pt}
\tracehead{\armGraphTools{}}
\tracestep{Modelo}{escolhe \texttt{query\_actions} e preenche os argumentos:}
\tracecode{filters: \{team: \jq Argentina\jq,\\
\ \ action: [\jq passe\jq], success: true,\\
\ \ third: \jq attacking\jq\},\\
group\_by: [\jq player\jq], top: 5}
\tracestep{Ferramenta}{executa a consulta fixa do seu código, com esses filtros, e devolve:}
\tracecode{Lionel Andrés Messi Cuccittini: 24\\ Alexis Mac Allister: 18 [...]}
\traceloop
\tracestep{Modelo}{responde: Lionel Andrés Messi Cuccittini}
\tracerule{0.8pt}
\captionof{figure}{Registro real de uma execução em cada abordagem com agente (\modelHaiku{}, pergunta a01),
lido na ferramenta de observabilidade. Em \armTextToCypher{}, o modelo escreve a consulta; em \armGraphTools{},
ele escolhe a ferramenta e preenche os argumentos, e a consulta já está no código. Nesta pergunta, uma chamada
bastou nas duas abordagens; em outras, o modelo fez várias. Saídas abreviadas com [...].}\label{fig:trace}
\end{minipage}\par\smallskip
